% SEGMENT OLP-0033-S01
% Part:sets-functions-relations
% Chapter: size-of-sets
% Section: non-enumerability

\documentclass[../../../include/open-logic-section]{subfiles}

\begin{document}

\olfileid{sfr}{siz}{nen}

\olsection{\printtoken{S}{nonenumerable} கணங்கள்}

\begin{editorial}
இந்தப் பிரிவு \olref[enm]{sec} இன் வரையறையைப் பயன்படுத்தி,
$\Bin^\omega$, $\Pow{\PosInt}$ ஆகியவற்றின் எண்ணிடத்தகாத
தன்மையை நிறுவுகிறது.
\olref[enm-alt]{sec} இல் உள்ள வடிவத்தை விட இது சற்றே
தொடக்கநிலையிலும் மேலும் விரிவாகவும் அமைக்கப்பட்டுள்ளது.
\end{editorial}

நேர்ம முழுக்களின் கணமான $\PosInt$ போன்ற சில கணங்கள் முடிவுறாதவை.
இதுவரை கண்ட முடிவுறாத கணங்களின் எடுத்துக்காட்டுகள் அனைத்தும்
!!{enumerable} ஆக இருந்தன. ஆனால் இந்தப் பண்பில்லாத முடிவுறாத
கணங்களும் உள்ளன. இத்தகைய கணங்கள்
\emph{!!{nonenumerable}} எனப்படும்.

முதலில், !!{nonenumerable} கணங்கள் இருக்கின்றன என்பதே
வியப்பாகத் தோன்றலாம்.
எந்த !!{enumerable} கணம் $A$ க்கும்
!!a{surjective} சார்பு $f \colon \PosInt \to A$ உள்ளது.
ஒரு கணம் !!{nonenumerable} எனில், இத்தகைய சார்பு இல்லை.
அதாவது, $\PosInt$ இன் முடிவுறாத எண்ணிக்கையிலான !!{element}s[acc]
$A$ க்கு அனுப்பும் எந்தச் சார்பாலும் $A$ முழுவதையும் அடைய முடியாது.
எனவே, முடிவுறாத எண்ணிக்கையிலான நேர்ம முழுக்களைவிட
$A$ இல் “மேலும்” !!{element}s உள்ளன.

ஒரு கணம் !!{nonenumerable} என்பதை எவ்வாறு நிறுவுவது?
இத்தகைய மேற்கோர்த்தல் சார்பு எதுவும் இருக்க முடியாது எனக் காட்ட வேண்டும்.
இதற்குச் சமானமாக, $A$ இன் உறுப்புகளை ஒரு திசையில் முடிவுறாமல்
செல்லும் பட்டியலில் பட்டியலாக்க முடியாது எனக் காட்ட வேண்டும்.
இதற்கான சிறந்த வழி, $A$ இன் !!{element}s கொண்ட ஒவ்வொரு பட்டியலும்
குறைந்தபட்சம் ஓர் உறுப்பை விட்டுவிட வேண்டும் என்று காட்டுவதாகும்;
அல்லது எந்தச் சார்பு $f\colon \PosInt \to A$ உம்
!!{surjective} பண்புடையதாக இருக்க முடியாது எனக் காட்டுவதாகும்.
காண்டோரின் \emph{மூலைவிட்ட முறையைப்} பயன்படுத்தி இதைச் செய்யலாம்.
$A$ இன் !!{element}s கொண்ட $x_1$, $x_2$, \dots என்ற
ஒரு பட்டியல் தரப்பட்டால், அதன் கட்டுமானத்தாலேயே அந்தப் பட்டியலில்
இடம்பெற முடியாத $A$ இன் மற்றோர் உறுப்பைக் கட்டமைக்கிறோம்.

நமது முதல் எடுத்துக்காட்டு, $0$, $1$ ஆகியவற்றின் எல்லா முடிவுறாத,
இடைவெளியற்ற தொடர்களையும் கொண்ட $\Bin^\omega$ என்ற கணமாகும்.

% SEGMENT OLP-0033-S02
\begin{thm}
\ollabel{thm:nonenum-bin-omega}
$\Bin^\omega$ என்பது !!{nonenumerable}.
\end{thm}

% SEGMENT OLP-0033-S03
\begin{proof}
முரண் பெறும் நோக்கில், $\Bin^\omega$ என்பது !!{enumerable} எனக் கொள்வோம்.
அதாவது, $\Bin^\omega$ இன் எல்லா !!{element}s கொண்ட
$s_{1}$, $s_{2}$, $s_{3}$, $s_{4}$, \dots{} என்ற பட்டியல்
ஒன்று உள்ளதாகக் கொள்வோம்.
இந்த $s_i$ ஒவ்வொன்றும் $0$, $1$ ஆகியவற்றின் முடிவுறாத தொடராகும்.
இந்தப் பட்டியலின் $i$ ஆவது தொடரின் $j$ ஆவது உறுப்பை
$s_i(j)$ என அழைப்போம்.
அப்போது $i$ ஆவது தொடர் $s_i$ என்பது:
\[
s_i(1), s_i(2), s_i(3), \dots
\]

இந்தப் பட்டியலையும், அதிலுள்ள ஒவ்வொரு தொடர் $s_i$ இன் உறுப்புகளையும்
ஓர் அட்டவணையில் அமைக்கலாம்:
\[
\begin{array}{c|c|c|c|c|c}
& 1 & 2 & 3 & 4 & \dots \\\hline
1 & \mathbf{s_{1}(1)} & s_{1}(2) & s_{1}(3) & s_1(4) & \dots \\\hline
2 & s_{2}(1)& \mathbf{s_{2}(2)} & s_2(3) & s_2(4) & \dots \\\hline
3 & s_{3}(1)& s_{3}(2) & \mathbf{s_3(3)} & s_3(4) & \dots \\\hline
4 & s_{4}(1)& s_{4}(2) & s_4(3) & \mathbf{s_4(4)} & \dots \\\hline
\vdots & \vdots & \vdots & \vdots & \vdots & \mathbf{\ddots}
\end{array}
\]
பக்கவாட்டில் உள்ள குறிகள் $s_1$, $s_2$, \dots என்ற பட்டியலிலுள்ள
தொடரின் எண்ணைக் காட்டுகின்றன.
மேற்புற எண்கள் ஒவ்வொரு தனித்த தொடரின் !!{element}s[acc]க் குறிக்கின்றன.
எடுத்துக்காட்டாக, $s_{1}(1)$ என்பது $s_{1}$ என்ற தொடரின்
முதல் !!{element} ஆக உள்ள எண், அதாவது $0$ அல்லது $1$, எதுவோ
அதற்கான பெயர்; இவ்வாறே தொடர்கிறது.

% SEGMENT OLP-0033-S04
இப்போது இந்தப் பட்டியலில் இடம்பெற முடியாத,
$0$, $1$ ஆகியவற்றின் ஒரு முடிவுறாத தொடர் $\overline{s}$ ஐக்
கட்டமைக்கிறோம். $\overline{s}$ இன் வரையறை
$s_1$, $s_2$, \dots என்ற பட்டியலைப் பொறுத்திருக்கும்.
$0$, $1$ ஆகியவற்றின் முடிவுறாத தொடர்களைக் கொண்ட எந்த முடிவுறாத
பட்டியலும், அந்தப் பட்டியலில் இடம்பெறாது என்பது உறுதியான
ஒரு முடிவுறாத தொடர் $\overline{s}$ ஐ உருவாக்குகிறது.

% SEGMENT OLP-0033-S05
$\overline{s}$ ஐ வரையறுக்க, அதன் எல்லா !!{element}s என்னவென்று,
அதாவது அனைத்து $n \in \PosInt$ க்கும் $\overline{s}(n)$ என்னவென்று,
குறிப்பிடுகிறோம்.
மேலுள்ள அட்டவணையின் மூலைவிட்டத்தில் கீழ்நோக்கிப் படித்து
(இதனாலேயே “மூலைவிட்ட முறை” என்ற பெயர் வருகிறது),
ஒவ்வொரு $1$ ஐயும் $0$ ஆகவும் ஒவ்வொரு $0$ ஐயும் $1$ ஆகவும்
மாற்றுவதன் மூலம் இதைச் செய்கிறோம்.
மேலும் அருவமாக, மூலைவிட்டத்தின் $n$ ஆவது !!{element} ஆன
$s_n(n)$ என்பது $1$ அல்லது $0$ ஆக இருப்பதைப் பொறுத்து,
$\overline{s}(n)$ ஐ $0$ அல்லது $1$ என வரையறுக்கிறோம்:
\[
\overline{s}(n) =
\begin{cases}
1 & \text{$s_{n}(n) = 0$ எனில்}\\
0 & \text{$s_{n}(n) = 1$ எனில்}.
\end{cases}
\]
வகைப்பிரிவு வரையறைகளைவிட வாய்பாடுகளை விரும்பினால்,
$\overline{s}(n) = 1 - s_n(n)$ எனவும் வரையறுக்கலாம்.

% SEGMENT OLP-0033-S06
$\overline{s}$ என்பது $0$, $1$ ஆகியவற்றின் ஒரு முடிவுறாத தொடர்
என்பது தெளிவு. ஏனெனில் அது நமது அட்டவணையின் மூலைவிட்டத்தில்
இடம்பெறும் $0$, $1$ ஆகியவற்றின் தொடரை எதிர்மாற்றிய தொடர் மட்டுமே.
எனவே $\overline{s}$ என்பது $\Bin^\omega$ இன் !!a{element}.
ஆனால் அது $s_1$, $s_2$, \dots{} என்ற பட்டியலில் இடம்பெற முடியாது.
ஏன்?

% SEGMENT OLP-0033-S07
அது பட்டியலின் முதல் தொடரான $s_1$ ஆக இருக்க முடியாது;
ஏனெனில் முதல் !!{element} இல் அது $s_1$ இலிருந்து வேறுபடுகிறது.
$s_1(1)$ எதுவாக இருந்தாலும், $\overline{s}(1)$ அதற்கு எதிரானதாக
இருக்குமாறு வரையறுத்தோம்.
அது பட்டியலின் இரண்டாவது தொடராகவும் இருக்க முடியாது;
ஏனெனில் இரண்டாவது உறுப்பில் $\overline{s}$ என்பது $s_2$ இலிருந்து
வேறுபடுகிறது: $s_2(2)$ என்பது $0$ எனில்,
$\overline{s}(2)$ என்பது $1$; மறுதிசையிலும் அவ்வாறே.
இவ்வாறே தொடர்கிறது.

% SEGMENT OLP-0033-S08
மேலும் துல்லியமாகக் கூறுவோம்:
$\overline{s}$ பட்டியலில் இருந்தால், ஏதாவது ஒரு $k$ க்கு
$\overline{s} = s_{k}$ என இருக்கும்.
இரண்டு தொடர்கள் ஒவ்வொரு இடத்திலும் ஒத்திருந்தால்,
அப்போதும் அப்போது மட்டுமே, அவை ஒன்றே.
அதாவது எந்த $n$ க்கும் $\overline{s}(n) =
s_{k}(n)$ ஆகும்.
குறிப்பாக $n = k$ என்ற சிறப்பு நிலையை எடுத்தால்,
$\overline{s}(k) = s_{k}(k)$ என இருக்க வேண்டும்.
$s_k(k)$ என்பது $0$ அல்லது $1$.
அது $0$ எனில், $\overline{s}(k)$ என்பது $1$ ஆக இருக்க வேண்டும்:
$\overline{s}$ ஐ அவ்வாறே வரையறுத்தோம்.
ஆனால் $s_k(k) = 1$ எனில், மீண்டும் $\overline{s}$ ஐ
நாம் வரையறுத்த விதத்தின்படி,
$\overline{s}(k) = 0$.
இரண்டு நிலைகளிலும் $\overline{s}(k) \neq s_{k}(k)$.

% SEGMENT OLP-0033-S09
$\Bin^\omega$ இன் !!{element}s கொண்ட
$s_1$, $s_2$, \dots{} என்ற பட்டியல் உள்ளது எனத் தொடக்கத்தில்
முன்கொண்டோம். இந்தப் பட்டியலிலிருந்து, பட்டியலில் இடம்பெற முடியாது
என நாம் நிறுவிய $\overline{s}$ என்ற தொடரைக் கட்டமைத்தோம்.
ஆனால் எல்லா $s_i$ களும் $0$, $1$ ஆகியவற்றின் தொடர்கள் எனில்,
அதுவும் நிச்சயமாக $0$, $1$ ஆகியவற்றின் தொடர்தான்;
அதாவது $\overline{s} \in
\Bin^\omega$.
குறிப்பாக, $\Bin^\omega$ இன் \emph{எல்லா} !!{element}s கொண்ட
பட்டியல் எதுவும் இருக்க முடியாது என்பதை இது காட்டுகிறது.
ஏனெனில் இத்தகைய எந்தப் பட்டியலுக்கும், அதில் இடம்பெறாது என்பது
உறுதியான $\overline{s}$ என்ற தொடரை மேலும் கட்டமைக்கலாம்.
ஆகவே $\Bin^\omega$ இன் எல்லாத் தொடர்களையும் கொண்ட
பட்டியல் இருக்கிறது என்ற முன்கோள் ஒரு முரணுக்கு வழிவகுக்கிறது.
\end{proof}

% SEGMENT OLP-0033-S10
\begin{explain}
இந்த நிறுவல் முறை அட்டவணையின் மூலைவிட்டத்தைப் பயன்படுத்தி
$\overline{s}$ ஐ வரையறுப்பதால் “மூலைவிட்டமாக்கம்” எனப்படுகிறது.
மூலைவிட்டமாக்கத்தில் ஓர் அட்டவணை இருக்க வேண்டும் என்பதில்லை.
அட்டவணையோ உண்மையான மூலைவிட்டமோ இல்லாதபோதும்,
ஒத்த கருத்தைப் பயன்படுத்திக் கணங்கள் !!{enumerable} அல்ல
எனக் காட்டலாம்.
\end{explain}

% SEGMENT OLP-0033-S11
\begin{thm}
\ollabel{thm:nonenum-pownat}
$\Pow{\PosInt}$ என்பது !!{enumerable} அல்ல.
\end{thm}

% SEGMENT OLP-0033-S12
\begin{proof}
$\PosInt$ இன் உட்கணங்களைக் கொண்ட ஒவ்வொரு பட்டியலுக்கும்,
அந்தப் பட்டியலில் இடம்பெற முடியாத $\PosInt$ இன் ஓர் உட்கணம்
உள்ளது எனக் காட்டி, அதே வழியில் செல்கிறோம்.
$\PosInt$ இன் உட்கணங்களைக் கொண்ட பின்வரும் பட்டியல்
தரப்பட்டுள்ளதாகக் கொள்வோம்:
\[
Z_{1}, Z_{2}, Z_{3}, \dots
\]
இப்போது எந்த $n \in \PosInt$ க்கும்,
$n \in \overline{Z}$ ஆக இருக்கும்போது, அப்போதும் அப்போது மட்டுமே,
$n \notin Z_{n}$ என அமையும் $\overline{Z}$ என்ற கணத்தை
வரையறுக்கிறோம்:
\[
\overline{Z} = \Setabs{n \in \PosInt}{n \notin Z_n}
\]
ஒவ்வொரு $Z_n$ உம் நேர்ம முழுக்களைக் கொண்ட கணம் என்பது முன்கோள்.
ஆகவே $\overline{Z}$ உம் நேர்ம முழுக்களைக் கொண்ட ஒரு கணம்;
இதனால் $\overline{Z} \in
\Pow{\PosInt}$.
ஆனால் $\overline{Z}$ பட்டியலில் இடம்பெற முடியாது.
இதைக் காட்ட, ஒவ்வொரு $k \in \PosInt$ க்கும்
$\overline{Z} \neq
Z_k$ என நிறுவுவோம்.

% SEGMENT OLP-0033-S13
எனவே ஏதேனும் ஒரு $k \in \PosInt$ ஐக் கொள்க.
எந்த $n \in \PosInt$ க்கும்,
$n \in \overline{Z}$ ஆக இருக்கும்போது, அப்போதும் அப்போது மட்டுமே,
$n \notin Z_n$ என அமையுமாறு $\overline{Z}$ ஐ வரையறுத்தோம்.
குறிப்பாக $n=k$ என எடுத்தால்,
$k \in \overline{Z}$ ஆக இருக்கும்போது, அப்போதும் அப்போது மட்டுமே,
$k \notin Z_k$.
ஆனால் $k$ என்பது இந்த இரு கணங்களில் ஒன்றின் !!a{element} ஆகவும்
மற்றொன்றின் உறுப்பு அல்லாததாகவும் இருப்பதால்,
$\overline{Z} \neq Z_k$ என்பதை இது காட்டுகிறது;
$\overline{Z}$, $Z_k$ ஆகியவற்றின் !!{element}s வெவ்வேறானவை.
$k$ ஏதேனும் ஒன்றாக இருந்ததால்,
$\overline{Z}$ என்பது $Z_1$, $Z_2$, \dots என்ற பட்டியலில் இல்லை.
\end{proof}

% SEGMENT OLP-0033-S14
\begin{explain}
முந்தைய நிறுவல் மூலைவிட்டத்தைக் குறிப்பிடவில்லை.
ஆனால் பின்வருமாறு காட்சிப்படுத்தினால், அதில் ஒரு மூலைவிட்டம்
இருப்பதாகக் கருதலாம்:
$Z_1$, $Z_2$, \dots என்ற கணங்கள் ஓர் அட்டவணையில் எழுதப்பட்டுள்ளதாகவும்,
ஒவ்வோர் !!{element} $j \in Z_i$ உம் $j$ ஆவது நெடுவரிசையில்
இடம்பெறுவதாகவும் கற்பனை செய்க.
அந்தப் பட்டியலின் முதல் நான்கு கணங்கள்
$\{1,2,3,\dots\}$, $\{2, 4, 6, \dots\}$,
$\{1,2,5\}$, $\{3,4,5,\dots\}$ எனக் கொள்வோம்.
அப்போது அட்டவணை பின்வருமாறு தொடங்கும்:
\[
\begin{array}{r@{}rrrrrrr}
  Z_1 = \{ & \mathbf{1}, & 2, & 3, & 4, & 5, & 6, & \dots\}\\
  Z_2 = \{ &  & \mathbf{2}, &  & 4, &  & 6, & \dots\}\\
  Z_3 = \{ & 1, & 2, &  &  & 5\phantom{,} &  & \}\\
  Z_4 = \{ &  &  & 3, & \mathbf{4}, & 5, & 6, & \dots\}\\
  \vdots & & & & & \ddots
\end{array}
\]
மூலைவிட்டத்தில் கீழ்நோக்கிச் சென்று, அங்கு இடம்பெறும் எண்களை
விட்டுவிட்டு, $j$ ஆவது கிடைவரிசை/நெடுவரிசையில் அட்டவணை
இடைவெளியாக உள்ள $j$ களைச் சேர்ப்பதால் கிடைக்கும் கணமே
$\overline{Z}$.
மேலுள்ள நிலையில், $1$, $2$ ஆகியவற்றை விட்டுவிட்டு,
$3$ ஐச் சேர்த்து, $4$ ஐ விட்டுவிடுவோம்; இவ்வாறே தொடர்கிறது.
\end{explain}

% SEGMENT OLP-0033-S15
\begin{prob}
ஒரு மூலைவிட்ட வாதத்தின் மூலம் $\Pow{\Nat}$ என்பது
!!{nonenumerable} எனக் காட்டுக.
\end{prob}

% SEGMENT OLP-0033-S16
\begin{prob}\label{sfr:siz:nen:prob:f-posint}
$f \colon \PosInt \to \PosInt$ என அமையும் எல்லாச் சார்புகளையும்
கொண்ட கணம் !!{nonenumerable} என்பதை ஒரு வெளிப்படையான
மூலைவிட்ட வாதத்தின் மூலம் காட்டுக.
அதாவது, $f_1$, $f_2$, \dots என்பது சார்புகளின் ஒரு பட்டியலாகவும்,
ஒவ்வொரு $f_i\colon
\PosInt \to \PosInt$ ஆகவும் இருந்தால்,
இந்தப் பட்டியலில் இல்லாத ஏதாவது ஒரு
$\overline{f}\colon \PosInt \to
\PosInt$ உள்ளது எனக் காட்டுக.
\end{prob}

\end{document}
